\documentclass{ximera}
\author{Jeffrey Kuan}
\title{Math 1050 HW 1.4 --- Exponents and the Order of Operations Agreement}
\license{CC: 0}

\begin{document}

\begin{abstract}
    Section 1.4 --- Exponents and the Order of Operations Agreement.
    Assigned exercises: 1–61 odd.
\end{abstract}

\maketitle

\begin{problem}
\textbf{1.} nine to the fifth power

\textbf{Answer:} $\answer{9^5}$
\end{problem}

\begin{problem}
\textbf{3.} seven to the $n$th power

\textbf{Answer:} $\answer{7^n}$
\end{problem}

\begin{problem}
\textbf{5.} $(-5)^2$, $-5^2$, and $-(5)^2$ all represent the same number.

\begin{multipleChoice}
\choice{True}
\choice[correct]{False}
\end{multipleChoice}
\end{problem}

\begin{problem}
\textbf{7.} To evaluate the expression $6 + 7 \cdot 10$ means to determine what one number it
  is equal to.

\begin{multipleChoice}
\choice{False}
\choice[correct]{True}
\end{multipleChoice}
\end{problem}

\begin{problem}
\textbf{9.} In the expression $(-5)^2$, $-5$ is called the \blank{} and $2$ is called the
  \blank. To evaluate $(-5)^2$, find the product $(\underline{\hspace{2.2em}})(\underline{\hspace{2.2em}}) = \underline{\hspace{2.2em}}$.

\textbf{Answer:} $\answer{\text{base}}$ and $\answer{\text{exponent}}$
\end{problem}

\begin{problem}
\textbf{11.} \emph{Evaluate.}\quad $6^2$

\textbf{Answer:} $\answer{36}$
\end{problem}

\begin{problem}
\textbf{13.} \emph{Evaluate.}\quad $-7^2$

\textbf{Answer:} $\answer{-49}$
\end{problem}

\begin{problem}
\textbf{15.} \emph{Evaluate.}\quad $(-3)^2$

\textbf{Answer:} $\answer{9}$
\end{problem}

\begin{problem}
\textbf{17.} \emph{Evaluate.}\quad $(-3)^4$

\textbf{Answer:} $\answer{81}$
\end{problem}

\begin{problem}
\textbf{19.} \emph{Evaluate.}\quad $\left(\dfrac{1}{2}\right)^{2}$

\textbf{Answer:} $\answer{\frac{1}{4}}$
\end{problem}

\begin{problem}
\textbf{21.} \emph{Evaluate.}\quad $(0.3)^2$

\textbf{Answer:} $\answer{0.09}$
\end{problem}

\begin{problem}
\textbf{23.} \emph{Evaluate.}\quad $2^2 \cdot (-3)$

\textbf{Answer:} $\answer{-12}$
\end{problem}

\begin{problem}
\textbf{25.} \emph{Evaluate.}\quad $2^3 \cdot 3^3 \cdot (-4)$

\textbf{Answer:} $\answer{-864}$
\end{problem}

\begin{problem}
\textbf{27.} \emph{Evaluate.}\quad $\left(\dfrac{2}{3}\right)^{2} \cdot 3^3$

\textbf{Answer:} $\answer{12}$
\end{problem}

\begin{problem}
\textbf{29.} \emph{Evaluate.}\quad $(0.3)^3 \cdot 2^3$

\textbf{Answer:} $\answer{0.216}$
\end{problem}

\begin{problem}
\textbf{31.} \emph{Evaluate.}\quad $\left(\dfrac{2}{3}\right)^{2} \cdot \dfrac{1}{4} \cdot 3^3$

\textbf{Answer:} $\answer{3}$
\end{problem}

\begin{problem}
\textbf{33.} Is the fifth power of negative eighteen positive or negative?

\begin{multipleChoice}
\choice[correct]{Negative}
\choice{Positive}
\end{multipleChoice}
\end{problem}

\begin{problem}
\textbf{35.} Why do we need an Order of Operations Agreement?

\begin{hint}
\textbf{Solution.} Without an agreed order, the same expression could be read in more than one way and give more than one value. For instance $2+3\cdot 4$ is $14$ if the multiplication is done first but $20$ if the addition is done first. The Order of Operations Agreement fixes one reading, so everyone who simplifies an expression gets the same result.
\end{hint}
\end{problem}

\begin{problem}
\textbf{37.} \emph{Complete Exercise 33 without actually finding the product.}\quad Simplify: $2\left(3^3\right)$

\textbf{Answer:} $\answer{54}$
\end{problem}

\begin{problem}
\textbf{39.} \emph{Simplify by using the Order of Operations Agreement.}\quad $4 - 8 \div 2$

\textbf{Answer:} $\answer{0}$
\end{problem}

\begin{problem}
\textbf{41.} \emph{Simplify by using the Order of Operations Agreement.}\quad $2(3 - 4) - (-3)^2$

\textbf{Answer:} $\answer{-11}$
\end{problem}

\begin{problem}
\textbf{43.} \emph{Simplify by using the Order of Operations Agreement.}\quad $24 - 18 \div 3 + 2$

\textbf{Answer:} $\answer{20}$
\end{problem}

\begin{problem}
\textbf{45.} \emph{Simplify by using the Order of Operations Agreement.}\quad $16 + 15 \div (-5) - 2$

\textbf{Answer:} $\answer{11}$
\end{problem}

\begin{problem}
\textbf{47.} \emph{Simplify by using the Order of Operations Agreement.}\quad $3 - 2\bigl[8 - (3 - 2)\bigr]$

\textbf{Answer:} $\answer{-11}$
\end{problem}

\begin{problem}
\textbf{49.} \emph{Simplify by using the Order of Operations Agreement.}\quad $6 + \dfrac{16 - 4}{2^2 + 2} - 2$

\textbf{Answer:} $\answer{6}$
\end{problem}

\begin{problem}
\textbf{51.} \emph{Simplify by using the Order of Operations Agreement.}\quad $96 \div 2\bigl[12 + (6 - 2)\bigr] - 3^3$

\textbf{Answer:} $\answer{-24}$
\end{problem}

\begin{problem}
\textbf{53.} \emph{Simplify by using the Order of Operations Agreement.}\quad $16 \div 2 - 4^2 - (-3)^2$

\textbf{Answer:} $\answer{-17}$
\end{problem}

\begin{problem}
\textbf{55.} \emph{Simplify by using the Order of Operations Agreement.}\quad $16 - 3(8 - 3)^2 \div 5$

\textbf{Answer:} $\answer{1}$
\end{problem}

\begin{problem}
\textbf{57.} \emph{Simplify by using the Order of Operations Agreement.}\quad $\dfrac{(-10) + (-2)}{6^2 - 30} \div |2 - 4|$

\textbf{Answer:} $\answer{-1}$
\end{problem}

\begin{problem}
\textbf{59.} \emph{Simplify by using the Order of Operations Agreement.}\quad $0.3(1.7 - 4.8) + (1.2)^2$

\textbf{Answer:} $\answer{0.51}$
\end{problem}

\begin{problem}
\textbf{61.} \emph{Simplify by using the Order of Operations Agreement.}\quad $\dfrac{3}{8} \div \left|\dfrac{5}{6} + \dfrac{2}{3}\right|$

\textbf{Answer:} $\answer{\frac{1}{4}}$
\end{problem}

\end{document}
